\documentclass[11pt]{article}
\usepackage[margin=0.85in]{geometry}
\usepackage{microtype}
\usepackage{parskip}

\title{Decision-Level vs. Cognition-Level Social Influence}
\author{Aslan}
\date{}

\begin{document}
\maketitle

Hi Sibi,

I now have two model configurations that use the same individual cognition,
decision, DeGroot influence, and trust-update components.  The difference is
where social influence is connected in the update sequence.

\textbf{Decision-level influence (the previous design).}
Everyone observes the same situation and updates their own cognitive state.
Each person then converts that state into a scalar preference for using
automation.  DeGroot influence is applied to these preferences, and the
driver's influenced preference determines the vehicle mode.

\textbf{Cognition-level influence (the alternative design).}
Everyone again observes the same situation and privately updates their
cognitive state.  DeGroot influence is then applied separately to automation
trust, perceived risk, and workload.  Each occupant forms an individual
preference and hypothetical decision from the influenced cognition.  Only the
driver's decision is executed; the passengers' decisions are retained for the
social-trust update and analysis.

Keeping these as two separate model classes lets me compare the connection
order while reusing the same components and parameters.

\section*{Preliminary result}

The simulations suggest that the models diverge most when occupants have
different \emph{internal cognitive profiles}, particularly when they favor the
same action for different reasons.  For example, the driver may favor
automation because perceived risk is low, while a passenger favors it because
trust in automation is high.  Their final preferences can therefore agree even
though their underlying cognition is very different.

In the synthetic high-contrast example, the decision-level model selected
automation in all eight cycles, whereas the cognition-level model selected
manual driving in all eight.  In a second example, where both occupants had
similar cognitive profiles and favored automation for the same reason, the two
models selected the same mode in every cycle; the mean absolute difference in
the driver's preference was only 0.0033.

In simple terms, discussing only final preferences can hide disagreement about
\emph{why} people prefer an action.  Discussing the cognitive components exposes
that disagreement, so ``influence first, then decide'' can differ substantially
from ``decide first, then influence.''  This occurs because the current
decision mapping takes the strongest of the trust, risk, and workload signals;
mixing cognitive states before this nonlinear operation is not generally
equivalent to mixing the resulting preferences afterward.

This is a result of the current synthetic comparison, not yet a general or
empirically validated conclusion.  A more precise working hypothesis is:

\begin{quote}
The difference between cognition-level and decision-level influence can become
large when occupants rely on different cognitive factors, especially when the
driver's result is close to the decision threshold.
\end{quote}

I think the cognition-level configuration is behaviorally plausible because
conversation can change a person's perceived risk, workload, or trust before a
decision is formed.  The decision-level configuration remains useful as the
original baseline.  My next step is to test this hypothesis systematically
with the participant-specific cognitive models rather than relying on these
illustrative parameter settings.

\end{document}
